\chapter*{\centering \fontsize{14pt}{16pt}\selectfont{\khmerfontmoul អត្ថបទសង្ខេប}}  
\addcontentsline{toc}{chapter}{\fontsize{12pt}{14pt}\selectfont{\khmerfontmoul អត្ថបទសង្ខេប}}

\begin{khmer}
\justifying
\khmerfontsr
\fontsize{12pt}{18pt}\selectfont
សន្ទស្សន៍ថ្លៃទំនិញប្រើប្រាស់ (\textenglish{CPI}) គឺជាសូចនាករគន្លឹះក្នុងការវាស់វែងអតិផរណា និងតាមដានកម្លាំងទិញជាក់ស្តែងរបស់ប្រជាពលរដ្ឋ ប៉ុន្តែការស្ទង់មតិបែបប្រពៃណីនៅកម្ពុជាមានគម្លាតពេលវេលាយឺតយ៉ាវពី ៣០ ទៅ ៦០ ថ្ងៃ និងចំណាយថវិកាច្រើន។ ដើម្បីដោះស្រាយបញ្ហានេះ គម្រោងស្រាវជ្រាវនេះបានបង្កើតប្រព័ន្ធវាស់វែងសន្ទស្សន៍ថ្លៃទំនិញប្រើប្រាស់ប្រចាំថ្ងៃនៅកម្ពុជា (\textenglish{\textbf{Cambodia Daily CPI Medallion Pipeline}}) ដោយប្រមូលទិន្នន័យថ្លៃទំនិញស្វ័យប្រវត្តជារៀងរាល់ព្រឹកនៅម៉ោង \textbf{០៨:០០ ព្រឹក} (\textenglish{ICT}) ពី ២៥ ប្រភពអនឡាញ គ្របដណ្តប់លើទាំង ១២ ជំពូក \textenglish{UN COICOP} ក្នុងរាជធានីភ្នំពេញ តាមរយៈស្ថាបត្យកម្ម \textenglish{PostgreSQL 16 Medallion (Bronze, Silver, Gold)}។

នៅក្នុងស្រទាប់ \textenglish{Silver} ប្រព័ន្ធអនុវត្តការផ្គូផ្គងអត្តសញ្ញាណទំនិញតាមរយៈ \textenglish{Vector Embeddings (\texttt{gemini-embedding-2})} និងវិធានការពារលក្ខណៈបច្ចេកទេសចំនួន ៧។ ចំណែកការចាត់ថ្នាក់ទំនិញចូលជំពូក ៤ ខ្ទង់ (\textenglish{DD.G.C}) ត្រូវបានអនុវត្តតាមយន្តការស្នូលពីរ៖ ហាងឯកទេសតែមួយត្រូវបានចាត់ថ្នាក់ដោយស្វ័យប្រវត្តតាមប្រព័ន្ធអាទិភាពហាង (\textenglish{\textbf{Store Priority Routing}}) ជាមួយល្បឿន $O(1)$ រីឯទំនិញថ្មីៗក្នុងផ្សារទំនើបចម្រុះត្រូវបានចាត់ថ្នាក់ជាបាច់ដោយបញ្ញាសិប្បនិម្មិត \textenglish{\textbf{Google Gemini AI (\texttt{gemini-2.5-flash})}} រួមជាមួយការចងចាំទិន្នន័យក្នុង \textenglish{PostgreSQL Memoization} និងវិធានការពារដែនកំណត់ទំនិញ (\textenglish{\textbf{Domain Trap Guardrails}}) ដើម្បីលុបបំបាត់កំហុសចាត់ថ្នាក់ឆ្លងជំពូក។

នៅក្នុងស្រទាប់ \textenglish{Gold} សន្ទស្សន៍ថ្លៃមូលដ្ឋានត្រូវបានគណនាតាមរូបមន្តមធ្យមធរណីមាត្រ \textenglish{\textbf{Jevons Formula}} ដើម្បីលុបបំបាត់លម្អៀងរូបមន្តគណិតវិទ្យា ($+1,18\%$/ឆ្នាំ) រួមជាមួយការបូកសរុបទម្ងន់ជាតិ \textenglish{CSES 2023}។ ម៉ាស៊ីនព្យាករណ៍អតិផរណាទាន់ពេល \textenglish{\textbf{Two-Stage Hybrid Ridge Nowcaster}} សម្រេចបានកម្រិតលម្អៀងជាក់ស្តែងត្រឹមតែ \textbf{+០,០១៨១ ភាគរយ} (+០,៣៦៥១\% ធៀបនឹង +០,៣៤៧០\% នៃទិន្នន័យផ្លូវការវិទ្យាស្ថានជាតិស្ថិតិ)។ លទ្ធផលទាំងអស់ត្រូវបានបង្ហាញផ្ទាល់តាមផ្ទាំងគ្រប់គ្រង \textenglish{Metabase} ចំនួន ៣ និង \textenglish{Power BI} សម្រាប់បម្រើការតាក់តែងគោលនយោបាយសេដ្ឋកិច្ចជាតិ។

\vspace{0.5cm}
\noindent \textbf{\khmerfont ពាក្យគន្លឹះ៖} សន្ទស្សន៍ថ្លៃទំនិញប្រើប្រាស់ (\textenglish{CPI}), ការប្រមូលទិន្នន័យស្វ័យប្រវត្ត (\textenglish{Web Scraping}), ស្ថាបត្យកម្ម \textenglish{Medallion}, ចំណាត់ថ្នាក់ \textenglish{COICOP}, បញ្ញាសិប្បនិម្មិត \textenglish{Gemini AI}, រូបមន្ត \textenglish{Jevons}, ការព្យាករណ៍អតិផរណា \textenglish{Ridge Nowcasting}។
\end{khmer}
\setmainfont{Times New Roman}